\chapter{Tecnologías}
\label{chap:tecnologias}

\lettrine{E}{n} este capítulo se describen las tecnologías software y hardware utilizadas en
el desarrollo del proyecto. Se detallan el lenguaje de programación empleado, el middleware
de comunicación que permite crear la arquitectura distribuida, las principales librerías y
herramientas integradas, así como los dispositivos físicos que permiten la interacción entre 
el sistema y los coches del circuito.


\section{Herramientas de desarollo software}

\subsection{Lenguajes de programacion}
\subsubsection{Python}

Python es un lenguaje de programación de código abierto, interpretado y de tipado dinámico, lo que significa que no requiere especificar el tipo de las variables de forma explícita y que el código se ejecuta directamente mediante un intérprete, sin necesidad de compilación previa~\cite{python}.

Admite múltiples paradigmas de programación, como la orientación a objetos, la programación funcional y la programación por procedimientos, y es multiplataforma, lo que permite ejecutar el mismo programa en distintos sistemas operativos sin modificaciones. Su extenso ecosistema de librerías científicas y de visión artificial, junto con su sintaxis clara, lo han convertido en uno de los lenguajes más populares en la actualidad.  En este proyecto, Python es el único lenguaje empleado en el sistema principal: todos los nodos de la arquitectura distribuida están implementados en Python sobre la librería \texttt{rclpy} de ROS2.
% TODO: Si documentas el firmware del Arduino, añade aquí un subsubsection breve sobre
% C/C++ (lenguaje del sketch), como hacen Mario y Adrián.
\subsubsection{C++}

\subsection{Librerías}
\subsubsection{OpenCV}
OpenCV es una librería de visión por computador de código abierto ampliamente utilizada~\cite{opencv,bradski2008}. En este proyecto es el componente central del análisis visual: se emplea para capturar imagenes desde las cámaras, convertir los fotogramas al espacio de color \gls{hsv}, segmentar las pegatinas de los coches y la línea de meta mediante máscaras de color, aplicar operaciones morfológicas de limpieza y extraer los contornos de los objetos detectados.
\subsubsection{NumPy}
NumPy es una librería optimizada para el cálculo numérico con arreglos multidimensionales~\cite{numpy}. Se utiliza para representar y operar sobre las trayectorias del circuito (arrays de puntos), calculando de forma vectorizada distancias entre puntos y segmentos, el nodo de la trayectoria más cercano a una posición o las máscaras binarias derivadas de los fotogramas. Estas operaciones se ejecutan decenas de veces por segundo, por lo que su eficiencia resulta esencial para el funcionamiento en tiempo real.

\subsubsection{pyserial}
Pyserial permite establecer comunicación serie entre Python y dispositivos externos como microcontroladores~\cite{pyserial}. En este proyecto se emplea en el nodo puente para enviar los comandos de velocidad al Arduino a través del puerto \gls{usb}, así como para recibir las confirmaciones de la placa.

\subsubsection{threading y concurrent.futures}
Los módulos estándar \texttt{threading} y \texttt{concurrent.futures} permiten ejecutar tareas en paralelo dentro de un mismo proceso. En el nodo cámara se emplean para desacoplar la captura de vídeo del procesamiento (un hilo dedicado actualizar el buffer de la cámara) y para buscar en paralelo las pegatinas de varios coches sobre el mismo fotograma mediante un \textit{pool} de hilos.

\subsection{Herramientas de despliegue}

\subsubsection{Docker y Docker Compose}

Docker es una plataforma de contenedores que permite empaquetar una aplicación junto con todas sus dependencias en una imagen ejecutable de forma idéntica en cualquier máquina~\cite{docker}. En este proyecto resulta clave para el despliegue distribuido: los equipos que participan en el sistema no necesitan tener instalado ROS2 ni ninguna librería, únicamente Docker. Una única imagen, construida sobre la base oficial \texttt{ros:humble-ros-base}, sirve para todos los roles del sistema.

Docker Compose complementa a Docker permitiendo describir en un fichero \gls{yaml} la configuración completa de arranque de un contenedor. Se definen dos ficheros: uno para las máquinas con cámara, parametrizado con el identificador del nodo y el dispositivo de vídeo, y otro para la máquina con acceso al puerto serie del Arduino. De esta forma, poner en marcha cualquier equipo del sistema se reduce a un único comando.

\subsection{Entornos de desarrollo}

\subsubsection{Visual Studio Code}

Visual Studio Code es un editor de código fuente gratuito y multiplataforma desarrollado por
Microsoft~\cite{vscode}. Ofrece resaltado de sintaxis, autocompletado, depuración integrada
y un amplio ecosistema de extensiones, entre ellas las de Python y ROS2, que facilitan el
desarrollo y la depuración del proyecto.

\subsubsection{Arduino IDE}

Es el entorno de desarrollo oficial utilizado para programar placas Arduino~\cite{arduinoide}.
Permite escribir el código del microcontrolador en un lenguaje basado en C/C++, compilarlo y
cargarlo en la placa mediante la conexión serie. En este proyecto se emplea para cargar el
firmware que interpreta los comandos de velocidad y genera las señales \gls{pwm} de los
carriles.

\subsection{Control de versiones}
\subsubsection{Git}
Git es un sistema de control de versiones distribuido~\cite{git}. Todo el desarrollo del
proyecto se ha gestionado con Git, empleando ramas para las líneas de trabajo principales y
una convención de mensajes de \textit{commit} con etiquetas por funcionalidad (por ejemplo
\texttt{[FUN-FINDCOLOR]}, o \texttt{[ALGO-CTRL]}) que facilita reconstruir la
evolución del sistema.

\subsection{Formatos de datos}

\subsubsection{YAML}

El formato \gls{yaml} permite representar datos estructurados de forma legible para humanos
y máquinas~\cite{pyyaml}. Es el formato estándar de configuración en ROS2 y de Docker
Compose. En este proyecto centraliza toda la configuración del sistema en un único fichero
de parámetros: colores de las pegatinas y de la meta, resolución de las cámaras, umbrales
del algoritmo de control o el puerto del Arduino, entre otros.

\subsubsection{JSON}

El formato \gls{json} es un estándar ligero de intercambio de datos basado en texto. Se
emplea para persistir en disco las trayectorias obtenidas durante la vuelta de calibración
(tanto las de cada cámara como las del controlador), de forma que en arranques posteriores
el sistema pueda cargarlas y omitir la calibración.

\section{ROS2}
\label{sec:ros2}

ROS2 (\textit{Robot Operating System 2}) es un conjunto de librerías y herramientas de código abierto para construir aplicaciones para robots~\cite{ros2docs}. A pesar de su nombre, no se trata de un sistema operativo, sino de un middleware que proporciona la infraestructura de comunicación, el sistema de construcción y las utilidades necesarias para componer aplicaciones a partir de procesos independientes llamados \textit{nodos}. En este trabajo se ha utilizado la distribución \textit{Humble Hawksbill}~\cite{ros2docs}, versión con soporte extendido (\gls{lts}) publicada en 2022.

ROS2 constituye la columna vertebral de la arquitectura de este proyecto, del que se emplean los siguientes conceptos. De cada uno se da aquí la idea necesaria para seguir el resto de la memoria, junto con la referencia a la documentación oficial en la que puede ampliarse.

\subsection{Nodos}
\label{sec:ros2-nodos}
Procesos independientes que encapsulan una funcionalidad concreta (en nuestro caso, la gestión de una cámara, el control de un coche o el puente con el Arduino) y que pueden ejecutarse en máquinas distintas~\cite{ros2nodes}.

\subsection{Topics y mensajes}
\label{sec:ros2-topics}
Bus que permite la comunicación entre los nodos que publican y reciben mensajes tipados~\cite{ros2topics}. En el proyecto se definen tipos de mensajes propios para intercambiar la posición del coche, el valor del pwm ("velocidad") y telemetría.

\subsection{Interfaces y mensajes propios}
\label{sec:ros2-interfaces}
Además de los tipos estándar, una aplicación puede definir los suyos en ficheros de texto con extensión \texttt{.msg}, en los que se declara un campo por línea con su tipo y su nombre~\cite{ros2interfaces}. Durante la compilación del paquete, ROS2 genera automáticamente a partir de ellos las clases que los nodos usan para publicar y recibir. Este proyecto define sus propios mensajes porque la información que intercambia (posiciones de dos pegatinas, consignas de PWM por carril, telemetría del lazo de control) no se corresponde con ningún tipo estándar.

\subsection{Parámetros}
\label{sec:ros2-parametros}
Valores de configuración declarados por cada nodo, que pueden establecerse desde un fichero \gls{yaml} común y modificarse en caliente durante la ejecución~\cite{ros2params}. Cada nodo puede además registrar una función de validación que se ejecuta antes de aceptar un cambio.

\subsection{Ficheros de lanzamiento}
\label{sec:ros2-launch}
Scripts de lanzamiento que permiten arrancar varios nodos junto con sus configuraciones~\cite{ros2launch}. Al describirse mediante código, la lista de nodos a levantar puede calcularse en el momento del arranque a partir de la configuración. Permiten también marcar un nodo para su reinicio automático (\textit{respawn}) si termina de forma inesperada.

\subsection{Namespaces}
\label{sec:ros2-namespaces}
Es un mecanismo de encapsulación jerárquica que agrupa identificadores del sistema (nodos, tópicos, servicios o parámetros) bajo un prefijo de ruta común, de forma análoga a los directorios de un sistema de ficheros~\cite{ros2names}. Un nodo lanzado dentro de un \textit{namespace} ve resueltos sus nombres relativos bajo ese prefijo, de modo que el mismo código sirve para varias instancias sin colisionar entre ellas.

\subsection{Executors y \emph{callback groups}}
\label{sec:ros2-executors}
Un nodo atiende su trabajo mediante \textit{callbacks}, funciones que se disparan al llegar un mensaje o al vencer un temporizador (\textit{timer}), una tarea periódica que el propio nodo programa. El \textit{executor} es quien decide qué \textit{callback} se ejecuta y en qué hilo~\cite{ros2executors}. El más sencillo los atiende de uno en uno; el \texttt{MultiThreadedExecutor} dispone de varios hilos y puede atenderlos a la vez. Los \textit{callback groups} controlan ese paralelismo: los de un grupo mutuamente excluyente (\texttt{MutuallyExclusiveCallbackGroup}) nunca se ejecutan simultáneamente entre sí, mientras que los de grupos distintos sí pueden hacerlo.

\subsection{Grabación de datos: rosbag2}
\label{sec:ros2-rosbag}
Utilidad que ROS2 incluye de serie para registrar en disco todos los mensajes que circulan por los \textit{topics} de un dominio y reproducirlos después~\cite{rosbag2}. El fichero resultante se denomina \textit{bag}. El formato de almacenamiento por omisión es mcap~\cite{mcap}, que guarda dentro del propio fichero la definición de cada tipo de mensaje registrado, de modo que una grabación sigue siendo legible aunque los mensajes cambien después.

\subsection{Colcon}
\label{sec:ros2-colcon}
Colcon es la herramienta oficial de construcción de "workspaces" de ROS2~\cite{colcon}.  Se encarga de compilar el paquete del proyecto, generar el código de los mensajes personalizados a partir de sus definiciones e instalar los ejecutables y ficheros de configuración para que puedan ser lanzados con las utilidades de ROS2.

\subsection{Paradigma publicador/suscriptor}
\label{sec:pubsub}

En el modelo publicador/suscriptor, los procesos no se comunican directamente entre sí, sino a través de canales con nombre (\textit{topics})~\cite{ros2topics}: los publicadores envían mensajes a un topic sin conocer quién los recibirá, y los suscriptores expresan su interés en un topic sin conocer quién publica en él. Este desacoplamiento aporta grandes ventajas que simplifican el diseño.

El paradigma encaja de forma natural en el problema de este proyecto: cada nodo cámara publica las posiciones que detecta sin saber cuántos controladores existen, y cada controlador se suscribe a las posiciones de su coche sin saber cuántas cámaras cubren el circuito. Añadir una cámara o un coche no requiere reconfigurar ningún otro componente del sistema.

\subsection{Middleware DDS y calidad de servicio (QoS)}
\label{sec:dds-qos}

\gls{dds} es un estándar de middleware publicador/suscriptor orientado a sistemas de tiempo real, definido por el Object Management Group~\cite{ddsspec}, sobre el que ROS2 construye toda su capa de comunicación. Dos de sus características son especialmente relevantes para este trabajo. La primera es el descubrimiento automático: los participantes de un mismo dominio (identificado por un número, el \textit{domain ID}~\cite{ros2domainid}) se anuncian periódicamente en la red local y establecen las conexiones entre publicadores y suscriptores sin ninguna configuración manual de direcciones y sin ningún proceso maestro. La segunda es la calidad de servicio (\gls{qos}~\cite{ros2qos}): cada publicador y suscriptor declara sus garantías de entrega mediante un conjunto de políticas, entre las que destacan:

\begin{itemize}
\item \textbf{Fiabilidad (\textit{reliability}):} en modo \textit{reliable} el middleware reenvía los mensajes perdidos hasta confirmar su entrega; en modo \textit{best effort} los mensajes se envían una única vez y pueden perderse.

\item \textbf{Durabilidad (\textit{durability}):} en modo \textit{volatile} un suscriptor solo recibe los mensajes publicados después de conectarse; en modo \textit{transient local} el publicador conserva los últimos mensajes y los entrega a los suscriptores que se conecten más tarde.

\item \textbf{Historial (\textit{history} y \textit{depth}):} determina cuántos mensajes se almacenan en las colas de envío y recepción.
\end{itemize}

Estas políticas se negocian entre las dos partes: un suscriptor solo recibe mensajes de un publicador si el perfil que declara es compatible con el de este, de modo que la elección del publicador condiciona a todos sus consumidores. ROS2 ofrece además perfiles predefinidos que agrupan valores habituales, como el orientado a datos de sensores~\cite{ros2qos}.

La elección de estas políticas supone un compromiso entre latencia y garantías de entrega. Para los datos de posición, que se publican decenas de veces por segundo, se emplea el perfil de datos de sensor (\textit{best effort}): un mensaje perdido carece de importancia porque enseguida llega el siguiente, y resulta preferible perder un fotograma a acumular retardo reenviando información ya obsoleta. En cambio, para la posición de la línea de meta, que se publica una única vez al arrancar, se utiliza un perfil \textit{reliable} y \textit{transient local}: el mensaje debe llegar con garantías incluso a los controladores que arranquen después de que la cámara lo haya publicado.

%\subsection{Tolerancia a fallos: heartbeat, failover y watchdog}
%\label{sec:tolerancia-fallos}
%
%En un sistema distribuido, la caída de un proceso no puede detectarse de forma directa: la
%única evidencia disponible es la ausencia de comunicación. Sobre esta idea se construyen los
%tres mecanismos de tolerancia a fallos empleados en este proyecto:
%
%\begin{itemize}
%    \item \textbf{Latido de vida (\textit{heartbeat}):} un proceso emite periódicamente un
%    mensaje que atestigua que sigue vivo. Quien lo vigila considera que ha fallado cuando
%    transcurre un plazo sin recibir latidos.
%
%    \item \textbf{Replicación primario/respaldo (\textit{failover}):} un componente crítico
%    se despliega por duplicado. La réplica primaria realiza el trabajo y emite latidos; la
%    réplica de respaldo procesa la misma información en silencio y, si detecta la ausencia
%    de latidos, asume el rol de primaria de forma automática. De este modo el servicio se
%    recupera en segundos sin intervención humana.
%
%    \item \textbf{Perro guardián (\textit{watchdog}):} un componente vigila que otro siga
%    enviando información con la cadencia esperada y, en caso contrario, ejecuta una acción
%    de seguridad. En este proyecto, el nodo puente detiene la alimentación de un carril
%    cuando deja de recibir consignas de velocidad para él, de forma que un coche cuyo
%    control ha fallado (o que se ha salido de la pista) queda detenido sin afectar al resto
%    de la carrera.
%\end{itemize}
%
%Como tercera capa de recuperación, el sistema de lanzamiento de ROS2 permite reiniciar
%automáticamente (\textit{respawn}) cualquier nodo que termine de forma inesperada, de manera
%que un nodo caído se reincorpora al sistema al cabo de unos segundos.

\section{Tecnologías hardware}
\label{sec:tecnoloxias-hardware}

\subsection{Circuito de Scalextric}

% TODO: Confirma marca y modelo del circuito y de los coches (Mario y Adrián usaron pista
% analógica Carrera línea EVOLUTION y coche Carrera 27495).
Se emplea un circuito analógico de coches tipo \textit{slot} con dos carriles alimentados de
forma independiente, de modo que la tensión aplicada a cada carril determina la velocidad
del coche que circula por él. Sobre los coches se colocan dos pegatinas de colores
distintivos, una delantera y otra trasera, que el sistema de visión utiliza para determinar
su posición y orientación. Además, sobre la pista se sitúa una marca de color en la línea de
meta, que el sistema detecta automáticamente para contabilizar vueltas y finalizar la
calibración. Este entorno físico constituye el escenario sobre el que interactúan la red de
cámaras y el control hardware del sistema.

\subsection{Placa Elegoo Uno R3}

La placa Elegoo Uno R3 es una réplica compatible de la placa Arduino Uno R3, basada en el
microcontrolador ATmega328P. Dispone de entradas y salidas digitales y analógicas y de
conexión \gls{usb}, y es totalmente compatible con el \gls{ide} oficial de Arduino. Su
función en el sistema es recibir por el puerto serie los comandos de velocidad enviados por
el nodo puente y generar las señales \gls{pwm} correspondientes para cada carril, actuando
como intermediario entre el software y el mundo físico.

\subsection{Módulo controlador L293D}

El L293D es un circuito integrado que implementa dos puentes en H, lo que permite controlar
dos flujos de corriente simultáneos e independientes. Se emplea habitualmente como
controlador de motores de corriente continua. 

\subsection{Webcams Microsoft LifeCam HD-5000}

% TODO: Confirma el modelo y el número de cámaras empleadas.
Las cámaras utilizadas son webcams Microsoft LifeCam HD-5000, capaces de capturar vídeo en
alta definición (720p) a 30 fotogramas por segundo a través de \gls{usb}. Cada cámara se
sitúa sobre una parte del trazado, de forma que el conjunto cubra el circuito completo. Para
reducir el ancho de banda \gls{usb} se solicita a la cámara el flujo comprimido en
\gls{mjpg}, y se desactiva el autoenfoque para evitar variaciones de imagen durante la
carrera.

\subsection{Raspberry Pi 5}
Estan equipadas con un procesador ARM Cortex-A76 de cuatro núcleos de 64 bits a 2,4 GHz, 4GB de RAM y puertos USB 3.0. En la arquitectura del sistema, este dispositivo se uso para conectar las cámaras y ejecutar los nodos de captura de ROS 2, adquiriendo las imágenes y publicando la posición del coche en tiempo real en la red para su posterior procesamiento.

\subsection{Raspberry Pi 4}
De manera complementaria, se emplean también placas Raspberry Pi 4 Modelo B, provistas de un procesador ARM Cortex-A72 de cuatro núcleos a 1,5 GHz, 4GB de RAM y puertos de comunicación USB. Cumplen una funcionalidad idéntica y paralela a la Raspberry Pi 5 dentro del proyecto.

\subsection{ASUS Vivobook S 14}
Esta equipado con un procesador AMD Ryzen AI 9 HX 370 de 12 núcleos a una frecuencia de hasta 5,16 GHz, gráficos integrados AMD Radeon 890M, 32 GB de memoria RAM LPDDR5X y 1 TB de almacenamiento SSD NVMe. El equipo opera bajo el sistema operativo Linux CachyOS en arquitectura x86\_64. Duranta las ejecuciones era donde iban los nodos de control y el nodo puente, al tener la conexión serie hacia el Arduino.